\lecture{7.2}{Augmented IPTW Estimation}
\label{lec:7-2}

\subsection{Combining G-Computation and IPTW}

G-computation and IPTW use different nuisance functions. G-computation predicts
counterfactual outcomes using $\Qbar_n(a,w)$. IPTW reweights observed outcomes
using $g_{n,a}(w)$. Augmented inverse probability weighting combines both ideas in a
single estimator.

The target counterfactual mean is
\[
\psi(a)=\E\{\Ya{a}\}.
\]
Under the usual identification assumptions,
\[
\psi(a)=\E\{\Qbar(a,W)\}
=
\E\left\{\frac{\ind(A=a)}{g_a(W)}Y\right\}.
\]
The AIPTW estimator adds a residual correction to the g-computation estimator, or
equivalently adds an outcome-regression correction to the IPTW estimator.

\subsection{AIPTW Estimator for $\E\{\Yone\}$}

For treatment level $1$, the AIPTW estimator is
\[
\psi_{n,\AIPTW}(1)
=
\frac{1}{n}\sum_{i=1}^n
\left[
\Qbar_n(1,W_i)
+
\frac{\ind(A_i=1)}{g_{n,1}(W_i)}
\{Y_i-\Qbar_n(1,W_i)\}
\right].
\]
The estimated ATE is obtained by forming the analogous estimator for treatment level
$0$ and subtracting:
\[
\widehat{\ATE}_{\AIPTW}
=
\psi_{n,\AIPTW}(1)-\psi_{n,\AIPTW}(0).
\]

\begin{estimationbox}
For binary treatment,
\[
\psi_{n,\AIPTW}(0)
=
\frac{1}{n}\sum_{i=1}^n
\left[
\Qbar_n(0,W_i)
+
\frac{\ind(A_i=0)}{g_{n,0}(W_i)}
\{Y_i-\Qbar_n(0,W_i)\}
\right],
\]
where $g_{n,0}(W_i)=1-g_{n,1}(W_i)$.
\end{estimationbox}

\subsection{IPTW Term}

The estimator can be rearranged as
\[
\psi_{n,\AIPTW}(1)
=
\frac{1}{n}\sum_{i=1}^n
\frac{\ind(A_i=1)}{g_{n,1}(W_i)}Y_i
+
\frac{1}{n}\sum_{i=1}^n
\left\{
1-\frac{\ind(A_i=1)}{g_{n,1}(W_i)}
\right\}
\Qbar_n(1,W_i).
\]
The first term is the usual IPTW estimator of $\E\{\Yone\}$:
\[
\frac{1}{n}\sum_{i=1}^n
\frac{\ind(A_i=1)}{g_{n,1}(W_i)}Y_i.
\]
It uses only observed treated outcomes, weighted by inverse treatment probabilities.

\subsection{Augmentation Term}

The second term,
\[
\frac{1}{n}\sum_{i=1}^n
\left\{
1-\frac{\ind(A_i=1)}{g_{n,1}(W_i)}
\right\}
\Qbar_n(1,W_i),
\]
is the augmentation term. It uses the outcome regression to correct the IPTW
estimator. If the estimated propensity score underweights or overweights certain
treated individuals, the augmentation term offsets part of the resulting distortion.

For example, suppose $Y$ is nonnegative and $g_{n,1}(W)$ tends to underestimate the
true probability of treatment. Then treated individuals receive weights that are too
large, and the IPTW part tends to overshoot. The factor
$1-\ind(A=1)/g_{n,1}(W)$ tends to be negative for treated observations with inflated
weights, so the augmentation term can pull the estimate downward.

\subsection{AIPTW as Augmented G-Computation}

The same estimator can be viewed as g-computation plus a weighted residual average:
\[
\psi_{n,\AIPTW}(1)
=
\underbrace{
\frac{1}{n}\sum_{i=1}^n \Qbar_n(1,W_i)
}_{\text{g-computation}}
+
\underbrace{
\frac{1}{n}\sum_{i=1}^n
\frac{\ind(A_i=1)}{g_{n,1}(W_i)}
\{Y_i-\Qbar_n(1,W_i)\}
}_{\text{augmentation}}.
\]
The augmentation is an average of observed residuals among treated individuals,
weighted to represent the target population. If the outcome regression
$\Qbar_n(1,W)$ systematically overpredicts treated outcomes, then the residuals
$Y-\Qbar_n(1,W)$ tend to be negative, and the augmentation term lowers the
g-computation estimate.

\subsection{Intuition for Bias Correction}

AIPTW uses $\Qbar_n$ to predict missing counterfactual outcomes and uses $g_n$ to
check those predictions against observed residuals in the treatment group that
actually reveals the relevant potential outcome. This creates a bias-correction
mechanism.

\begin{figure}[H]
\centering
\begin{tikzpicture}[
  scale=.86,
  transform shape,
  node distance=.85cm,
  box/.style={draw, rounded corners=2pt, minimum width=2.65cm, minimum height=.76cm},
  arr/.style={-{Latex[length=2.5mm]}, thick}
]
\node[box] (gcomp) {Predict everyone\\ $\Qbar_n(1,W_i)$};
\node[box, right=of gcomp] (resid) {Compute treated\\ residuals};
\node[box, right=of resid] (weight) {Weight residuals\\ by $1/g_{n,1}(W_i)$};
\node[box, right=of weight] (add) {Add correction\\ to mean};
\draw[arr] (gcomp) -- (resid);
\draw[arr] (resid) -- (weight);
\draw[arr] (weight) -- (add);
\end{tikzpicture}
\caption{AIPTW augments predicted counterfactual means with weighted residual corrections.}
\label{fig:ch7-aiptw-correction}
\end{figure}

\begin{keyidea}
AIPTW is both augmented IPTW and augmented g-computation. The two algebraic forms
are identical, but they highlight different sources of protection: the propensity score
helps correct outcome-regression errors, and the outcome regression helps correct
propensity-score errors.
\end{keyidea}
